\documentclass[reqno]{amsart}
\usepackage[T1]{fontenc}
\usepackage{amsmath,amssymb,amsthm}
\usepackage{xr-hyper}
\usepackage[hidelinks]{hyperref}
\externaldocument[tt-]{20261125T004100--tate-twists-and-sha__workbook}
\externaldocument[ec-]{20261125T014500--eisenstein-congruences-and-selmer-classes__workbook}
\externaldocument[dp-]{20261126T030900--denominators-of-projections-to-lifts__workbook}

% Chapter number; \ref{chap} gives it.
\newcounter{chap}
% Chapter numbers, generated from 20261125T012000--grh-workbook__guide.org
% by scripts/make-chapters.py.  Do not edit: change the order of the
% index and rerun the script (./preview-site.sh and CI do this).
%
% Usage in a sheet, after \newcounter{chap}:
%   % Chapter numbers, generated from 20261125T012000--grh-workbook__guide.org
% by scripts/make-chapters.py.  Do not edit: change the order of the
% index and rerun the script (./preview-site.sh and CI do this).
%
% Usage in a sheet, after \newcounter{chap}:
%   % Chapter numbers, generated from 20261125T012000--grh-workbook__guide.org
% by scripts/make-chapters.py.  Do not edit: change the order of the
% index and rerun the script (./preview-site.sh and CI do this).
%
% Usage in a sheet, after \newcounter{chap}:
%   \input{chapters}\chapnumber{<sheet ID>}
\makeatletter
\providecommand\chapters@set[2]{\@namedef{chapters@#1}{#2}}
\chapters@set{20261125T001000}{1}% zeros-to-ramp-sums
\chapters@set{20261125T001100}{2}% inert-primes-and-estimate-b
\chapters@set{20261125T005000}{3}% galois-fourier-and-averaging
\chapters@set{20261125T005500}{4}% columns-from-a-sphere
\chapters@set{20261125T000000}{5}% slope-method
\chapters@set{20261126T030100}{6}% a-twist-for-each-prime
\chapters@set{20261125T003000}{7}% the-chain-and-the-square-root
\chapters@set{20261125T010000}{8}% rank-lemma-and-zero-estimate
\chapters@set{20261125T010500}{9}% poincare-series-and-near-orthogonality
\chapters@set{20261124T223000}{10}% elliptic-curves-for-the-proof
\chapters@set{20261125T009500}{11}% fock-model-and-theta-intertwiners
\chapters@set{20261124T223100}{12}% complex-multiplication-and-inert-reduction
\chapters@set{20261125T008000}{13}% one-cm-determinant
\chapters@set{20261124T214500}{14}% supersingular-theta-divisibility
\chapters@set{20261124T214600}{15}% the-isotropic-plane
\chapters@set{20261125T013700}{16}% pairings-on-a-formal-group-of-height-two
\chapters@set{20261125T013500}{17}% hecke-nondegeneracy-on-the-tree
\chapters@set{20261126T030200}{18}% the-bound-is-attained
\chapters@set{20261125T009000}{19}% kronecker-limit-and-chowla-selberg
\chapters@set{20261125T002000}{20}% definite-theta-lifts
\chapters@set{20261125T015500}{21}% siegel-weil-in-the-smallest-case
\chapters@set{20261125T002100}{22}% definite-theta-lifts-and-norms
\chapters@set{20261125T016500}{23}% hecke-operators-and-theta-lifts
\chapters@set{20261125T016700}{24}% divisibility-along-the-chain
\chapters@set{20261125T004000}{25}% denominators-and-sha
\chapters@set{20261125T004100}{26}% tate-twists-and-sha
\chapters@set{20261126T020430}{27}% the-minus-part-by-gauss-sums
\chapters@set{20261125T014100}{28}% congruence-modules
\chapters@set{20261125T014500}{29}% eisenstein-congruences-and-selmer-classes
\chapters@set{20261125T014700}{30}% pseudocharacters-and-the-length-bound
\chapters@set{20261125T011000}{31}% galois-representations-black-box
\chapters@set{20261126T030300}{32}% the-data
\chapters@set{20261126T020210}{33}% bessel-poincare-series-and-their-gram-operators
\chapters@set{20261126T030400}{34}% the-seesaw-identity
\chapters@set{20261126T030500}{35}% divisibility-of-the-composite-lift
\chapters@set{20261126T030600}{36}% theorem-a-short-forms
\chapters@set{20261126T030700}{37}% the-determinant-of-a-chain
\chapters@set{20261126T030800}{38}% the-cost-of-one-lift
\chapters@set{20261126T030900}{39}% denominators-of-projections-to-lifts
\chapters@set{20261126T031000}{40}% twists-with-few-congruences
\chapters@set{20261126T031100}{41}% theorem-c-and-the-main-inequality
\providecommand\chapnumber[1]{%
  \@ifundefined{chapters@#1}%
    {\PackageError{chapters}{Sheet #1 is not listed in the index}%
      {List it in 20261125T012000--grh-workbook__guide.org\MessageBreak
       and run scripts/make-chapters.py.}}%
    {\edef\thechap{\@nameuse{chapters@#1}}}}
\makeatother
\chapnumber{<sheet ID>}
\makeatletter
\providecommand\chapters@set[2]{\@namedef{chapters@#1}{#2}}
\chapters@set{20261125T001000}{1}% zeros-to-ramp-sums
\chapters@set{20261125T001100}{2}% inert-primes-and-estimate-b
\chapters@set{20261125T005000}{3}% galois-fourier-and-averaging
\chapters@set{20261125T005500}{4}% columns-from-a-sphere
\chapters@set{20261125T000000}{5}% slope-method
\chapters@set{20261126T030100}{6}% a-twist-for-each-prime
\chapters@set{20261125T003000}{7}% the-chain-and-the-square-root
\chapters@set{20261125T010000}{8}% rank-lemma-and-zero-estimate
\chapters@set{20261125T010500}{9}% poincare-series-and-near-orthogonality
\chapters@set{20261124T223000}{10}% elliptic-curves-for-the-proof
\chapters@set{20261125T009500}{11}% fock-model-and-theta-intertwiners
\chapters@set{20261124T223100}{12}% complex-multiplication-and-inert-reduction
\chapters@set{20261125T008000}{13}% one-cm-determinant
\chapters@set{20261124T214500}{14}% supersingular-theta-divisibility
\chapters@set{20261124T214600}{15}% the-isotropic-plane
\chapters@set{20261125T013700}{16}% pairings-on-a-formal-group-of-height-two
\chapters@set{20261125T013500}{17}% hecke-nondegeneracy-on-the-tree
\chapters@set{20261126T030200}{18}% the-bound-is-attained
\chapters@set{20261125T009000}{19}% kronecker-limit-and-chowla-selberg
\chapters@set{20261125T002000}{20}% definite-theta-lifts
\chapters@set{20261125T015500}{21}% siegel-weil-in-the-smallest-case
\chapters@set{20261125T002100}{22}% definite-theta-lifts-and-norms
\chapters@set{20261125T016500}{23}% hecke-operators-and-theta-lifts
\chapters@set{20261125T016700}{24}% divisibility-along-the-chain
\chapters@set{20261125T004000}{25}% denominators-and-sha
\chapters@set{20261125T004100}{26}% tate-twists-and-sha
\chapters@set{20261126T020430}{27}% the-minus-part-by-gauss-sums
\chapters@set{20261125T014100}{28}% congruence-modules
\chapters@set{20261125T014500}{29}% eisenstein-congruences-and-selmer-classes
\chapters@set{20261125T014700}{30}% pseudocharacters-and-the-length-bound
\chapters@set{20261125T011000}{31}% galois-representations-black-box
\chapters@set{20261126T030300}{32}% the-data
\chapters@set{20261126T020210}{33}% bessel-poincare-series-and-their-gram-operators
\chapters@set{20261126T030400}{34}% the-seesaw-identity
\chapters@set{20261126T030500}{35}% divisibility-of-the-composite-lift
\chapters@set{20261126T030600}{36}% theorem-a-short-forms
\chapters@set{20261126T030700}{37}% the-determinant-of-a-chain
\chapters@set{20261126T030800}{38}% the-cost-of-one-lift
\chapters@set{20261126T030900}{39}% denominators-of-projections-to-lifts
\chapters@set{20261126T031000}{40}% twists-with-few-congruences
\chapters@set{20261126T031100}{41}% theorem-c-and-the-main-inequality
\providecommand\chapnumber[1]{%
  \@ifundefined{chapters@#1}%
    {\PackageError{chapters}{Sheet #1 is not listed in the index}%
      {List it in 20261125T012000--grh-workbook__guide.org\MessageBreak
       and run scripts/make-chapters.py.}}%
    {\edef\thechap{\@nameuse{chapters@#1}}}}
\makeatother
\chapnumber{<sheet ID>}
\makeatletter
\providecommand\chapters@set[2]{\@namedef{chapters@#1}{#2}}
\chapters@set{20261125T001000}{1}% zeros-to-ramp-sums
\chapters@set{20261125T001100}{2}% inert-primes-and-estimate-b
\chapters@set{20261125T005000}{3}% galois-fourier-and-averaging
\chapters@set{20261125T005500}{4}% columns-from-a-sphere
\chapters@set{20261125T000000}{5}% slope-method
\chapters@set{20261126T030100}{6}% a-twist-for-each-prime
\chapters@set{20261125T003000}{7}% the-chain-and-the-square-root
\chapters@set{20261125T010000}{8}% rank-lemma-and-zero-estimate
\chapters@set{20261125T010500}{9}% poincare-series-and-near-orthogonality
\chapters@set{20261124T223000}{10}% elliptic-curves-for-the-proof
\chapters@set{20261125T009500}{11}% fock-model-and-theta-intertwiners
\chapters@set{20261124T223100}{12}% complex-multiplication-and-inert-reduction
\chapters@set{20261125T008000}{13}% one-cm-determinant
\chapters@set{20261124T214500}{14}% supersingular-theta-divisibility
\chapters@set{20261124T214600}{15}% the-isotropic-plane
\chapters@set{20261125T013700}{16}% pairings-on-a-formal-group-of-height-two
\chapters@set{20261125T013500}{17}% hecke-nondegeneracy-on-the-tree
\chapters@set{20261126T030200}{18}% the-bound-is-attained
\chapters@set{20261125T009000}{19}% kronecker-limit-and-chowla-selberg
\chapters@set{20261125T002000}{20}% definite-theta-lifts
\chapters@set{20261125T015500}{21}% siegel-weil-in-the-smallest-case
\chapters@set{20261125T002100}{22}% definite-theta-lifts-and-norms
\chapters@set{20261125T016500}{23}% hecke-operators-and-theta-lifts
\chapters@set{20261125T016700}{24}% divisibility-along-the-chain
\chapters@set{20261125T004000}{25}% denominators-and-sha
\chapters@set{20261125T004100}{26}% tate-twists-and-sha
\chapters@set{20261126T020430}{27}% the-minus-part-by-gauss-sums
\chapters@set{20261125T014100}{28}% congruence-modules
\chapters@set{20261125T014500}{29}% eisenstein-congruences-and-selmer-classes
\chapters@set{20261125T014700}{30}% pseudocharacters-and-the-length-bound
\chapters@set{20261125T011000}{31}% galois-representations-black-box
\chapters@set{20261126T030300}{32}% the-data
\chapters@set{20261126T020210}{33}% bessel-poincare-series-and-their-gram-operators
\chapters@set{20261126T030400}{34}% the-seesaw-identity
\chapters@set{20261126T030500}{35}% divisibility-of-the-composite-lift
\chapters@set{20261126T030600}{36}% theorem-a-short-forms
\chapters@set{20261126T030700}{37}% the-determinant-of-a-chain
\chapters@set{20261126T030800}{38}% the-cost-of-one-lift
\chapters@set{20261126T030900}{39}% denominators-of-projections-to-lifts
\chapters@set{20261126T031000}{40}% twists-with-few-congruences
\chapters@set{20261126T031100}{41}% theorem-c-and-the-main-inequality
\providecommand\chapnumber[1]{%
  \@ifundefined{chapters@#1}%
    {\PackageError{chapters}{Sheet #1 is not listed in the index}%
      {List it in 20261125T012000--grh-workbook__guide.org\MessageBreak
       and run scripts/make-chapters.py.}}%
    {\edef\thechap{\@nameuse{chapters@#1}}}}
\makeatother
\chapnumber{20261125T014700}
\renewcommand{\thesection}{\thechap.\arabic{section}}

\theoremstyle{definition}
\newtheorem{problem}{Problem}[section]
\theoremstyle{plain}
\newtheorem{theorem}{Theorem}[section]

\newcommand{\Z}{\mathbb{Z}}
\newcommand{\Q}{\mathbb{Q}}
\newcommand{\F}{\mathbb{F}}
\newcommand{\Zp}{\Z_p}
\newcommand{\Qp}{\Q_p}
\newcommand{\Oc}{\mathcal{O}}
\newcommand{\Tb}{\mathbf{T}}
\newcommand{\Frob}{\mathrm{Frob}}
\newcommand{\tr}{\operatorname{tr}}
\newcommand{\len}{\operatorname{length}}
\newcommand{\Sha}{\mathrm{Sha}}
\newcommand{\GQ}{G_{\Q,p}}
\newcommand{\Eis}{\mathrm{Eis}}
\newcommand{\mat}[4]{\bigl(\begin{smallmatrix} #1 & #2 \\ #3 & #4 \end{smallmatrix}\bigr)}

\setlength{\emergencystretch}{2em}
\title{Pseudocharacters and the length bound}
\author{}
\date{November 25, 2026}

\begin{document}
\maketitle
\refstepcounter{chap}\label{chap}

\paragraph{Prerequisites.}
Complete local noetherian rings and Hensel's lemma; the Artin--Rees lemma \cite[Cor.~10.10]{AM}; lengths of modules; elementary divisors of matrices over \(\Zp\); graded rings.  From Chapter~\ref{tt-chap}, \emph{Odd Tate twists}: the groups \(\GQ\), \(G_p\), \(\Sha_p(\Q, e)\), Problem~\ref{tt-pr:H0} and Theorem~\ref{tt-thm:MW}.  From Chapter~\ref{ec-chap}, \emph{Eisenstein congruences and Selmer classes}: Problem~\ref{ec-pr:691}, the congruence~\eqref{ec-eq:congruence}, the lattices of Section~\ref{ec-sec:lattices} (Problems~\ref{ec-pr:stable}--\ref{ec-pr:ends}), the inputs (D), (I), (O) of Section~\ref{ec-sec:ribet} with Problems~\ref{ec-pr:ss} and~\ref{ec-pr:unram}, Section~\ref{ec-sec:tangent}, in particular Problems~\ref{ec-pr:red-ideal}--\ref{ec-pr:691-segment}, and Section~\ref{ec-sec:shapiro}.

Section~\ref{ec-sec:tangent} of Chapter~\ref{ec-chap} observed that the Hecke algebra sees a Galois representation only through its traces.  The Hecke algebra on the Eisenstein and cuspidal eigenspaces in Problem~\ref{pr:hecke-691} has two components.  In general the traces in a Hecke algebra need not be realized by a representation over that algebra.  What structure do traces alone carry, and how do they measure \(\Sha\)?  The structure is a pseudocharacter; the ideal of products \(b \cdot c\) of Chapter~\ref{ec-chap} can be read off from the traces; and the quotient by it carries extension classes, which bound the congruence from above by the length of a Selmer group.  In the eigenform example, the cocycle obeys every linear relation among the matrices \(U(g)-1\) for the twist \(U=\chi^{1-k}\rho_f\) (Problem~\ref{pr:kappa-obeys}).  Theorem~\ref{thm:graph} realizes such cocycles by equivariant maps on submodules of families of matrices.  In a graded model of the deformation ring we finally compute the reverse comparison.  Throughout, \(p\) is an odd prime.

\section{Pseudocharacters of dimension two}\label{sec:pseudo}

Let \(G\) be a group and \(R\) a commutative ring in which \(2\) is invertible.  A \emph{pseudocharacter of dimension~\(2\)} of \(G\) with values in \(R\) is a pair \((T, D)\) of a function \(T \colon G \to R\) and a homomorphism \(D \colon G \to R^\times\) such that \(T(1) = 2\) and, for all \(g, h \in G\),
\[
  T(gh) = T(hg), \qquad T(gh) + D(g)\, T(g^{-1}h) = T(g)\, T(h) .
\]
By \cite[Lemma~1.9]{Ch} such pairs are the same as Chenevier's determinants of dimension~\(2\).  The function \(T\) is then a pseudocharacter of dimension~\(2\) in the sense of Taylor \cite[\S1]{Ta} and Rouquier \cite{Ro}, and conversely every such \(T\) arises in this way \cite[Prop.~1.29]{Ch}; see also \cite[\S1.2]{BC}.  The notion is due to Wiles in dimension two and to Taylor in general \cite[\S1.1]{BC}.

\begin{problem}\label{pr:trace-pc}
For every homomorphism \(\rho \colon G \to \mathrm{GL}_2(R)\), \((\tr\rho, \det\rho)\) is a pseudocharacter.  In particular, so is \((\psi_1 + \psi_2, \psi_1\psi_2)\) for any two characters \(\psi_1, \psi_2 \colon G \to R^\times\).
\end{problem}

\begin{problem}\label{pr:D-from-T}
If \((T, D)\) is a pseudocharacter, then \(D(g) = \tfrac12\bigl(T(g)^2 - T(g^2)\bigr)\) and \(T(g^{-1}) = D(g)^{-1} T(g)\) for all \(g \in G\).
\end{problem}

So \(D\) is determined by \(T\), and we often write \(T\) for the pair.

\begin{problem}\label{pr:subring}
Let \(R \subseteq R_1 \times R_2\) be a subring in which \(2\) is invertible, and let \((T_i, D_i)\) be pseudocharacters with values in \(R_i\).  If \((T_1(g), T_2(g)) \in R\) for every \(g\), then \(((T_1, T_2), (D_1, D_2))\) is a pseudocharacter with values in \(R\).
\end{problem}

The first example of a pseudocharacter that is not visibly a trace is the one carried by a Hecke algebra with one Eisenstein and one cuspidal point.  Let \(G\) be profinite, let \(\chi \colon G \to \Zp^\times\) be a continuous character whose reduction \(\bar\chi\) is not trivial, and let \(\rho \colon G \to \mathrm{GL}_2(\Qp)\) be a continuous representation with \(\tr\rho(G) \subseteq \Zp\), \(\tr \rho \equiv 1 + \chi \pmod p\) and \(\tr\rho \ne 1 + \chi\).  Put
\[
  e := \min_{g \in G} v_p\bigl(\tr\rho(g) - 1 - \chi(g)\bigr), \qquad \mathcal{T}(g) := \bigl(1 + \chi(g), \tr\rho(g)\bigr) \in \Zp \times \Zp .
\]

\begin{problem}\label{pr:toy-hecke}
The \(\Zp\)-subalgebra of \(\Zp \times \Zp\) generated by the \(\mathcal{T}(g)\) is \(\Tb_e := \{(u, v) : u \equiv v \pmod{p^e}\}\), and \(x \mapsto (0, p^e)\) induces \(\Zp[x]/(x(x - p^e)) \cong \Tb_e\).  The function \(\mathcal{T}\) is a pseudocharacter with values in \(\Tb_e\).
\end{problem}

The ring \(\Tb_e\) is the \emph{toy Hecke algebra}: \(x = 0\) is the Eisenstein point, \(x = p^e\) the cuspidal point, and \(e\) measures how close they are.  It occurs in nature.

\begin{problem}\label{pr:hecke-691}
The \(\Z_{691}\)-subalgebra of \(\mathrm{End}(M_{12}(\Q) \otimes \Q_{691})\) generated by the Hecke operators \(T_n\), \(n \ge 1\), on modular forms of weight \(12\) and level \(1\) is isomorphic to \(\Z_{691}[x]/(x(x - 691))\).
\end{problem}

\section{The reducibility ideal}\label{sec:red}

Let \(R\) be a complete local noetherian ring with maximal ideal \(\mathfrak m\) and residue field \(\F\) of characteristic \(\ne 2\) (so \(2 \in R^\times\)), and let \((T, D)\) be a pseudocharacter of \(G\) with values in \(R\) such that \(T \equiv \chi_1 + \chi_2 \pmod{\mathfrak m}\) for two distinct characters \(\chi_1, \chi_2 \colon G \to \F^\times\).  Fix \(\sigma_0 \in G\) with \(\chi_1(\sigma_0) \ne \chi_2(\sigma_0)\).

\begin{problem}\label{pr:roots}
The polynomial \(X^2 - T(\sigma_0)X + D(\sigma_0)\) has roots \(\lambda_1, \lambda_2 \in R\) with \(\lambda_i \equiv \chi_i(\sigma_0) \pmod{\mathfrak m}\), and \(\lambda_1 - \lambda_2 \in R^\times\).
\end{problem}

Define \(a, d \colon G \to R\) by
\[
  a(g) := \frac{T(\sigma_0 g) - \lambda_2 T(g)}{\lambda_1 - \lambda_2} , \qquad d(g) := T(g) - a(g) .
\]

\begin{problem}\label{pr:a-entry}
Suppose that \(T = \tr\rho\) for \(\rho \colon G \to \mathrm{GL}_2(R)\) with \(\rho(\sigma_0) = \mathrm{diag}(\lambda_1, \lambda_2)\), and write \(\rho = \mat{a'}{b}{c}{d'}\).  Then \(a = a'\), \(d = d'\), and \(a(gh) - a(g)a(h) = b(g)c(h)\).
\end{problem}

So \(a\) and \(d\) are the diagonal entries of a representation that need not exist, and the defect of \(a\) from being multiplicative stands in for the products \(b(g)c(h)\).

\begin{problem}\label{pr:a-residual}
\(a \equiv \chi_1\) and \(d \equiv \chi_2 \pmod{\mathfrak m}\).  In particular \(a(g) \in R^\times\).
\end{problem}

\begin{problem}\label{pr:D-a}
\(D(g) = a(g)T(g) - a(g^2)\) for every \(g \in G\).
\end{problem}

The \emph{reducibility ideal} \(I \subseteq R\) is the ideal generated by the elements \(a(gh) - a(g)a(h)\), \(g, h \in G\).

\begin{problem}\label{pr:mod-I}
Modulo \(I\), \(a\) and \(d\) are characters \(G \to (R/I)^\times\), \(T \equiv a + d\) and \(D \equiv ad\).
\end{problem}

\begin{problem}\label{pr:minimal}
Let \(J \subsetneq R\) be an ideal, and let \(\psi_1, \psi_2 \colon G \to (R/J)^\times\) be characters with \(\psi_i \equiv \chi_i \pmod{\mathfrak m}\) and \(T \equiv \psi_1 + \psi_2 \pmod J\).  Then \(\psi_1 \equiv a\) and \(\psi_2 \equiv d \pmod J\), and \(I \subseteq J\).
\end{problem}

Problems~\ref{pr:mod-I} and~\ref{pr:minimal} prove the following; in particular \(I\) does not depend on \(\sigma_0\).  Over a henselian local ring in which \(d!\) is invertible, Bellaïche and Chenevier give the corresponding reducibility ideal for each partition of the pairwise distinct, absolutely irreducible residual constituents of a \(d\)-dimensional pseudocharacter \cite[Prop.~1.5.1]{BC}.

\begin{theorem}\label{thm:red}
\(I\) is the smallest ideal \(J\) of \(R\) such that \(T\) is congruent modulo \(J\) to a sum of two characters lifting \(\chi_1\) and \(\chi_2\).
\end{theorem}

For a representation \(\rho\) over a discrete valuation ring as in Section~\ref{ec-sec:lattices}, Problems~\ref{pr:a-entry} and~\ref{ec-pr:red-ideal} give \(I = \pi^{m+n}\Oc\), where \(m + n\) is the length of the segment of stable lattices.  Return to the toy Hecke algebra.

\begin{problem}\label{pr:I-toy}
In Problem~\ref{pr:toy-hecke}, suppose that \(\rho\) is irreducible, and let \(m + n\) be the length of its segment of stable lattices (Problem~\ref{ec-pr:segment}, with \(\chi_1 = 1\), \(\chi_2 = \bar\chi\)).  Then \(m + n \ge e\), and the reducibility ideal of \(\mathcal{T}\) is \(\{0\} \times p^{m+n}\Zp = p^{m+n-e}\, x\Tb_e\).
\end{problem}

The \emph{Eisenstein ideal} of a pseudocharacter \(T\) with respect to two given characters \(\psi_1, \psi_2 \colon G \to R^\times\) is the ideal \(I_\Eis\) generated by the \(T(g) - \psi_1(g) - \psi_2(g)\).  By Theorem~\ref{thm:red}, \(I \subseteq I_\Eis\) whenever the \(\psi_i\) lift the \(\chi_i\).  In the setting of Chapter~\ref{ec-chap} the two ideals coincide.

\begin{problem}\label{pr:seg-eis}
Let \(k\), \(p\), \(f\), \(\lambda\) be as in Section~\ref{ec-sec:ribet}, and let \(m + n\) be the length of the segment of stable lattices of \(\rho_f\).  Then
\[
  m + n = \min_{\ell \ne p} v\bigl(a_\ell(f) - 1 - \ell^{k-1}\bigr) ,
\]
the minimum over primes \(\ell \ne p\).  So the reducibility ideal of \(\tr\rho_f\) equals its Eisenstein ideal for \((1, \chi^{k-1})\).
\end{problem}

For \(\Delta\) and \(p = 691\) both sides are \(1\) (Problem~\ref{ec-pr:691-segment}).  When \(K = \Qp\), in the toy Hecke algebra built from \(\rho = \rho_f\) and \(\chi^{k-1}\) we get \(e = m + n\), and the reducibility ideal is \(x\Tb_e\): the Eisenstein component \(\Tb_e/x\Tb_e = \Zp\) is the largest quotient on which \(\mathcal{T}\) is reducible.

\section{The quotient by the reducibility ideal and \texorpdfstring{\(\mathrm{Ext}^1\)}{Ext1}}\label{sec:ext}

Ribet's lemma produced a class modulo \(\pi\).  The reducibility ideal produces a class modulo \(I\), which generates a copy of \(\Oc/I\).  Keep the notation of Section~\ref{ec-sec:lattices}, with \(\chi_1 \ne \chi_2\), and write \(\rho\) in the basis \(\pi^m e_1, e_2\) of \(L_m\) as \(\mat{a}{b}{c}{d}\).  Then the values of \(b\) generate \(\Oc\), the values of \(c\) generate \(I = \pi^{m+n}\Oc\), and by Problem~\ref{pr:mod-I} the reductions \(\psi_1, \psi_2\) of \(a, d\) modulo \(I\) are characters.  Put \(\psi := \psi_1\psi_2^{-1}\).

\begin{problem}\label{pr:cocycle}
\(\kappa(g) := b(g)\,d(g)^{-1} \bmod I\) is a cocycle on \(G\) with values in \((\Oc/I)(\psi)\).
\end{problem}

\begin{problem}\label{pr:exact-order}
The \(\Oc\)-submodule of \(H^1(G, (\Oc/I)(\psi))\) generated by the class of \(\kappa\) is isomorphic to \(\Oc/I\).
\end{problem}

Thus the quotient \(\Oc/I\) embeds in a group of extensions of \(\psi_2\) by \(\psi_1\).  For two distinct residual characters over a henselian local ring \(R\), Bellaïche and Chenevier construct an injection from \(\operatorname{Hom}_R(B,R/I)\) to the algebraic extension group of \(\psi_2\) by \(\psi_1\), where \(B\) is the upper right module of a chosen Cayley--Hamilton quotient \cite[Thm.~1.5.5]{BC}; for the continuous extension group under additional hypotheses, see \cite[Prop.~1.5.10(ii)]{BC}.  Now take \(G = \GQ\), and \(f\), \(\lambda\), \(\Oc\) as in Problem~\ref{pr:seg-eis}, so that \(\psi \equiv \chi^{1-k} \pmod I\).  As in Section~\ref{ec-sec:shapiro}, \(\Sha_p(\Q, k-1) \otimes_{\Zp} \Oc\) is the group of classes in \(H^1(\GQ, (K/\Oc)(1-k))\) that vanish on \(G_p\).  Use the coefficient embedding \((\Oc/I)(1-k)\hookrightarrow(K/\Oc)(1-k)\), \(\bar a\mapsto a/\pi^{m+n}+\Oc\).

\begin{problem}\label{pr:sha-lower}
The class of \(\kappa\) vanishes on \(G_p\), and its image in \(H^1(\GQ, (K/\Oc)(1-k))\) generates a submodule isomorphic to \(\Oc/I\).  Hence \(\len_\Oc\bigl(\Sha_p(\Q, k-1) \otimes \Oc\bigr) \ge m + n\).
\end{problem}

This is the direction of Ribet and of Mazur and Wiles \cite{MW}, for one eigenform: a congruence of exponent \(m + n\) produces a cyclic piece of \(\Sha\) of that length.  For \(\Delta\) it gives \(\len \Sha_{691}(\Q, 11) \ge 1\), and Theorem~\ref{tt-thm:MW} gives equality.  Read the other way, Problem~\ref{pr:sha-lower} bounds a congruence by a Selmer group: \(m + n \le \len_\Oc\bigl(\Sha_p(\Q, k-1) \otimes \Oc\bigr)\).  In the rewritten Lean proof \cite{GL}, a Selmer length contributes, in this direction, to the upper bound for the cost of clearing a projection.  The corresponding statement of the rewritten Lean proof is \texttt{exists\_operator\_ellVal\_le\_liftValuationBound} (\texttt{New/Congruences/FirstOrder/DenominatorBound.lean}); see Chapter~\ref{dp-chap}, \emph{Denominators of projections to lifts}.  In Problem~\ref{pr:sha-lower} the class already lies in the kernel of restriction to \(G_p\), so its full length contributes to the lower bound for \(\Sha\).

\section{Cocycles that obey the relations}\label{sec:relations}

The module \texttt{ribetCocycles} of the rewritten Lean proof \cite{GL} consists of cocycles obeying the simultaneous relations below (\texttt{mem\_ribetCocycles}, \texttt{New/Congruences/LocalLoss/Basic.lean}).  Problem~\ref{pr:kappa-obeys} gives an example.  Let \(\Lambda\) be a commutative ring, \(G\) a group, \(X\) a finite set, and, for \(\iota \in X\), let \(U_\iota \colon G \to M_q(\Lambda)\) be homomorphisms of monoids; let \(\varrho \colon G \to M_s(\Lambda)\) be another, and \(\alpha \in \Lambda\).  Put \(\mathcal{N} := M_q(\Lambda)^X\), on which \(g\) acts by \((m_\iota)_\iota \mapsto (U_\iota(g)m_\iota)_\iota\), and \(U(g) - 1 := (U_\iota(g) - 1)_\iota \in \mathcal{N}\).  A map \(\gamma \colon G \to (\Lambda/\alpha)^s\) is a \emph{cocycle for \(\varrho\)} if \(\gamma(gh) = \gamma(g) + \varrho(g)\gamma(h)\).  It \emph{obeys the relations among the \(U(g) - 1\)} if, for every finitely supported \(\theta \colon G \to \Lambda\),
\[
  \sum_{g} \theta(g)\,\bigl(U(g) - 1\bigr) = 0 \quad\Longrightarrow\quad \sum_{g} \theta(g)\,\gamma(g) = 0 .
\]
If all \(U_\iota\) are trivial, only \(\gamma = 0\) obeys them.

\begin{problem}\label{pr:kappa-obeys}
In the setting of Problem~\ref{pr:sha-lower}, put \(\Lambda := \Oc\), \(\alpha := \pi^{m+n}\), \(U := \chi^{1-k}\rho_f\) in the basis \(\pi^m e_1, e_2\) of \(L_m\) (one representation, \(q = 2\)) and \(\varrho := \chi^{1-k}\) (\(s = 1\)).  Then \(\kappa(g) = \chi(g)^{1-k} b(g) \bmod I\), and \(\kappa\) is a cocycle for \(\varrho\) that obeys the relations among the \(U(g) - 1\).
\end{problem}

From now on let \(\gamma\) be a cocycle for \(\varrho\) that obeys the relations among the \(U(g) - 1\), and let \(\mathcal{M} \subseteq \mathcal{N}\) be the \(\Lambda\)-span of the \(U(g) - 1\), \(g \in G\).

\begin{problem}\label{pr:graph-map}
There is a unique \(\Lambda\)-linear map \(\phi \colon \mathcal{M} \to (\Lambda/\alpha)^s\) with \(\phi(U(g) - 1) = \gamma(g)\) for all \(g\).  The module \(\mathcal{M}\) is stable under \(G\), and \(\phi(g y) = \varrho(g)\phi(y)\) for \(y \in \mathcal{M}\).
\end{problem}

So the graph of \(\phi\) is a \(G\)-stable submodule of \(\mathcal{N} \times (\Lambda/\alpha)^s\), with \(g\) acting on the second factor through \(\varrho\), on which the first projection is injective.  It is usually too small: even over \(\Zp\), \(\mathcal{M}\) need not contain \(p^j\mathcal{N}\) for any \(j \ge 0\).  Let \(\Lambda\) be noetherian.  By the Artin--Rees lemma \cite[Cor.~10.10]{AM}, applied to the submodule \(\mathcal{M}\) of the finitely generated module \(\mathcal{N}\), there is \(\nu \ge 1\) with \(\mathcal{M} \cap \alpha^\nu \mathcal{N} \subseteq \alpha\mathcal{M}\).

\begin{problem}\label{pr:graph-AR}
Let \(\mathcal{M} \cap \alpha^\nu \mathcal{N} \subseteq \alpha\mathcal{M}\).  Then
\[
  Y := \bigl\{ \bigl(y + \alpha^\nu v,\ \phi(y)\bigr) : y \in \mathcal{M},\ v \in \mathcal{N} \bigr\}
\]
is a \(G\)-stable submodule of \(\mathcal{N} \times (\Lambda/\alpha)^s\), the first projection is injective on \(Y\), and \(\bigl(U(g) - 1 + \alpha^\nu v, \gamma(g)\bigr) \in Y\) for all \(g \in G\) and \(v \in \mathcal{N}\).
\end{problem}

\begin{problem}\label{pr:graph-toy}
Let \(\Lambda = \Zp\), \(\alpha = p\), \(G = \Z\), \(X\) a point, \(U(j) = \mat{1}{jp}{0}{1}\), \(\varrho = 1\) and \(\gamma(j) = j \bmod p\).  Then \(\gamma\) obeys the relations among the \(U(j) - 1\), and the least \(\nu\) with \(\mathcal{M} \cap p^\nu M_2(\Zp) \subseteq p\mathcal{M}\) is \(\nu = 2\).  For \(\nu = 1\) the set \(Y\) of Problem~\ref{pr:graph-AR} contains \((0, 0)\) and \((0, 1)\).
\end{problem}

Problems~\ref{pr:graph-map} and~\ref{pr:graph-AR} prove the following theorem.  This is \texttt{ribetCocycles.exists\_graph\_subrepresentation} in the rewritten Lean proof \cite{GL}.

\begin{theorem}\label{thm:graph}
Let \(\Lambda\) be a commutative noetherian ring, \(\alpha \in \Lambda\), \(G\) a group, \(X\) a finite set, and \(U_\iota \colon G \to M_q(\Lambda)\), \(\iota \in X\), and \(\varrho \colon G \to M_s(\Lambda)\) homomorphisms of monoids.  Let \(\gamma \colon G \to (\Lambda/\alpha)^s\) be a cocycle for \(\varrho\) that obeys the relations among the \(U(g) - 1\).  Then there are \(\nu \ge 0\) and a \(G\)-stable \(\Lambda\)-submodule \(Y\) of \(M_q(\Lambda)^X \times (\Lambda/\alpha)^s\), with \(g\) acting by \(\bigl((m_\iota)_\iota, w\bigr) \mapsto \bigl((U_\iota(g)m_\iota)_\iota, \varrho(g)w\bigr)\), such that the first projection is injective on \(Y\) and \(\bigl((U_\iota(g) - 1)_\iota + \alpha^\nu v, \gamma(g)\bigr) \in Y\) for all \(g \in G\) and all \(v \in M_q(\Lambda)^X\).
\end{theorem}

In the application over an \(\ell\)-adic coefficient ring with \(\alpha=\ell^m\), the first projection identifies \(Y\) with a lattice in the space of matrix families.  The proof of \texttt{comesFromClass\_of\_mem\_ribetCocycles} uses this lattice to describe the restriction of the cocycle.

\paragraph{In the Lean formalization.}
In the rewritten Lean proof \cite{GL}, Theorem~\ref{thm:graph} is \texttt{ribetCocycles.exists\_graph\_subrepresentation}, in \texttt{New/Congruences/LocalLoss/Basic.lean} of the readable edition, and the cocycles of its hypothesis form the module \texttt{ribetCocycles}.  The local and global applications below use a prime \(\ell\).  In \texttt{comesFromClass\_of\_mem\_ribetCocycles} (\texttt{New/Congruences/LocalLoss/TateDualBound.lean}), let \(\Oc\) be the ring of integers of a finite extension of \(\Q_\ell\), let \(\eta\colon\Gamma\to G\) be a homomorphism, and give \(\Gamma\) a homomorphism to the absolute Galois group of \(\Q_\ell\) whose image contains a neighborhood of the identity.  Fix a class \(\mathcal C=\texttt{New.Lift.InClass}\,F\,K'\,U'\) of lattices, specified by a filter \(F\) on \(\Gamma\), a predicate \(K'\) on group elements and polynomials over \(\Oc\), and a finite set \(U'(\lambda')\) of allowed Hodge--Tate jumps for each \(\Z_\ell\)-algebra map \(\lambda'\colon\Oc\to\mathbb C_\ell\).  If the lattices of all \(U_\iota\circ\eta\) belong to \(\mathcal C\), then for every \(m\ge0\) and every cocycle \(w\colon G\to(\Oc/\ell^m\Oc)^s\) for \(\varrho\) obeying the relations among the \(U(g)-1\), one has \(w\circ\eta=q\circ a\), where \(a\) is a crystalline cocycle of a lattice \(Y\) in \(\mathcal C\) and \(q\colon Y\to(\Oc/\ell^m\Oc)^s\) is \(\Oc\)-linear and equivariant for \(\varrho\circ\eta\).  Under the hypotheses of \texttt{exists\_length\_le\_loss\_add\_crystallineSelmerLength} (\texttt{New/Congruences/FirstOrder/RibetElement.lean}), let \(\Oc\) be its coefficient ring with fraction field \(K\), a finite Galois extension of \(\Q_\ell\) in \(\overline{\Q}_\ell\), and let \(T\) have rank \(r'\).  Write \(\mathrm{sel}_T\) for the \(\Oc\)-length of the Selmer module of locally constant cocycles for \(T\) in \((K/\Oc)^{r'}\) with the prescribed ramification and crystalline relations and with value zero at a specified element \(\gamma_0\).  There are \(b_0\in\mathbb N\) and \(\mathrm{loss}_T\in\mathbb N\cup\{\infty\}\) such that the cocycles for \(T\) in \((K/\Oc)^{r'}\) obeying the relations among the \(U(g)-1\) and vanishing at the chosen tame generators, the inverses of their complex conjugates, and \(\gamma_0\) form a module of length at most \(b_0=\mathrm{loss}_T+\mathrm{sel}_T\).  The proof takes \(\mathrm{loss}_T\) to be the sum, over the places above \(\ell\), of the local defects of the Tate dual of \(T\).  The proof of \texttt{ribet\_columns} (same file) uses this length bound to control a coefficient of Ribet's element, and that estimate enters \texttt{exists\_operator\_ellVal\_le\_liftValuationBound} and then \texttt{denominatorLe\_clearingBound}, the bound for the orthogonal projection (\texttt{New/Congruences/FirstOrder/DenominatorBound.lean}).

\section{A graded model and the length bound}\label{sec:model}

At the Eisenstein point the ordinary first-order lifts of Problem~\ref{ec-pr:trace-blind} all have the same trace.  A product of upper and lower entries has weight \(0\).  The reverse of the inequality of Problem~\ref{pr:sha-lower} would bound the length of \(\Sha\) by trace differences.  Conjugation by \(\mathrm{diag}(1,t)\) in Section~\ref{ec-sec:tangent} gives the upper right and lower left entries weights \(-1\) and \(+1\), respectively, and fixes traces.  We compare the two sides in a polynomial model with these weights, in which everything can be computed.  This section is a complement; later chapters do not use it.

Let \(P \in \Zp[u]\) be a polynomial with \(P(0) \ne 0\) and leading coefficient in \(\Zp^\times\), and let
\[
  A := \Zp[y, z]/\bigl(y\, P(yz)\bigr) ,
\]
graded by giving \(y\) weight \(-1\) and \(z\) weight \(+1\); the relation is homogeneous of weight \(-1\).  Think of \(y\) as an upper right entry and \(z\) as a lower left entry, so that \(u = yz\) plays the role of \(b \cdot c\).  Let \(x \colon A \to \Zp\) be the Eisenstein point \(y = z = 0\), with ideal \(I_x = (y, z)\).  The \emph{cotangent module} \(I_x/I_x^2\) is graded; write \((I_x/I_x^2)_{-1}\) for its part of weight \(-1\), the model for the dual of the Selmer group.

\begin{problem}\label{pr:model}
\(A_0 = \Zp[u]/(u P(u))\) with \(u = yz\), and \((I_x/I_x^2)_{-1} \cong \Zp/P(0)\Zp\).  For an integer \(r \ge 1\) and \(P(u) = p^r - u\), \(A_0\) is the toy Hecke algebra \(\Tb_r\) of Problem~\ref{pr:toy-hecke}.
\end{problem}

\begin{problem}\label{pr:points}
A point \((y, z) \in \overline{\Q}_p^2\) of \(\operatorname{Spec} A\) has \(y = 0\), or else \(yz\) is a root of \(P\) and \(z \ne 0\).  For each root \(u_0\) of \(P\), the points with \(yz = u_0\) form one orbit of the torus.
\end{problem}

Call a point with both entries nonzero \emph{two-signed}.  The torus orbits of two-signed \(\overline{\Q}_p\)-points correspond to the distinct roots of \(P\), or equivalently to the \(\overline{\Q}_p\)-points of the cuspidal quotient \(\Tb^0:=\Zp[u]/(P(u))\) of \(A_0\).  The ideal \(I_x\cap A_0=(u)\) has image \(I_\Eis:=u\Tb^0\) in this quotient; as in Section~\ref{sec:red}, its generator \(u = yz\) is the product of an upper and a lower entry.

\begin{problem}\label{pr:norm}
\(\Tb^0\) is free over \(\Zp\).  For \(a \in \Zp[u]\) with \(a(\xi) \ne 0\) at every root \(\xi\) of \(P\),
\[
  \len_{\Zp} \Tb^0/a\Tb^0 = \sum_\xi v(a(\xi)) ,
\]
the sum over the roots of \(P\) in \(\overline{\Q}_p\) with multiplicity, where \(v(p) = 1\).
\end{problem}

\begin{problem}\label{pr:length-bound}
Let \(a \in A_0\) vanish at \(x\), and let its image in \(\Tb^0\) be nonzero at every root of \(P\).  Then
\[
  \len (I_x/I_x^2)_{-1} = \len \Tb^0/I_\Eis \le \sum_\xi v(\xi(a)) ,
\]
where \(\xi(a)\) is the value of \(a\) at \(u = \xi\) and the sum is over the roots \(\xi\) of \(P\) in \(\overline{\Q}_p\), with multiplicity; that is, over the torus orbits of two-signed points (Problem~\ref{pr:points}), each counted with the multiplicity of its root.  Equality holds for \(a = u\).
\end{problem}

Problems~\ref{pr:model}--\ref{pr:length-bound} prove the following.

\begin{theorem}\label{thm:model}
In the model \(A\), the length of the cotangent module of weight \(-1\) equals \(\len \Tb^0/I_\Eis\), and is at most \(\sum_\xi v(\xi(a))\) for every \(a\) as in Problem~\ref{pr:length-bound}.  Equality holds for \(a = u\).
\end{theorem}

For \(p = 691\) and \(P(u) = 691 - u\) the model has \(A_0 = \Z_{691}[x]/(x(x-691))\), the Hecke algebra of Problem~\ref{pr:hecke-691}, one two-signed orbit, with \(u = 691\), and a cotangent module of weight \(-1\) of length \(1 = \len\Sha_{691}(\Q, 11)\).  In the model two things are built in: the cotangent module stands for the dual of the Selmer group, and every two-signed point is cuspidal.  For an actual deformation ring both must be proved; the second is an automorphy theorem for residually reducible representations (Section~\ref{ec-sec:tangent}).

\section*{Hints}

\subsection*{\S\ref{sec:pseudo}}
\begin{description}
\item[\ref{pr:trace-pc}] Cayley--Hamilton: \(\rho(g) + \det\rho(g)\,\rho(g)^{-1} = \tr\rho(g)\).
\item[\ref{pr:D-from-T}] Take \(h = g\), then \(h = 1\).
\item[\ref{pr:subring}] Problem~\ref{pr:D-from-T}.
\item[\ref{pr:toy-hecke}] Subtract the scalar \(1 + \chi(g)\) from \(\mathcal{T}(g)\).  Then Problems~\ref{pr:D-from-T}, \ref{pr:trace-pc} and~\ref{pr:subring}.
\item[\ref{pr:hecke-691}] Use the eigenbasis \(E_{12}, \Delta\), Problem~\ref{ec-pr:691}, and \(\tau(2) = -24\).
\end{description}

\subsection*{\S\ref{sec:red}}
\begin{description}
\item[\ref{pr:roots}] Hensel's lemma.
\item[\ref{pr:a-entry}] Compute \(\tr(\rho(\sigma_0)\rho(g))\).
\item[\ref{pr:a-residual}] Reduce modulo \(\mathfrak m\) and use \(T(\sigma_0 g) \equiv \chi_1(\sigma_0 g) + \chi_2(\sigma_0 g)\).
\item[\ref{pr:D-a}] Express \(T(\sigma_0 g^2)\) by the defining identity with \((g, g\sigma_0)\), and \(T(g^2)\) by Problem~\ref{pr:D-from-T}.
\item[\ref{pr:mod-I}] Problems~\ref{pr:a-residual} and~\ref{pr:D-a}.
\item[\ref{pr:minimal}] Problem~\ref{pr:D-from-T}, and the polynomial of Problem~\ref{pr:roots} modulo \(J\).
\item[\ref{pr:I-toy}] First, the reduction of a stable lattice has semisimplification \(1 \oplus \bar\chi\): its trace is \(1 + \bar\chi\), its determinant is \(\bar\chi\) by Problem~\ref{pr:D-from-T}, and Brauer--Nesbitt applies.  The roots of Problem~\ref{pr:roots} and the function \(a\) may be computed in each factor.  Use Problem~\ref{pr:a-entry} in the second factor, and note that \(\Tb_e\) projects onto the second factor.
\item[\ref{pr:seg-eis}] One inequality is Problem~\ref{ec-pr:trace-cong}.  For the other, apply Problem~\ref{pr:minimal} with \(J\) the ideal generated by the \(a_\ell(f) - 1 - \ell^{k-1}\), and Chebotarev.
\end{description}

\subsection*{\S\ref{sec:ext}}
\begin{description}
\item[\ref{pr:cocycle}] Modulo \(I\) the matrix is upper triangular; multiply two of them.
\item[\ref{pr:exact-order}] Evaluate at \(\sigma_0\).
\item[\ref{pr:sha-lower}] Run Problem~\ref{ec-pr:unram} modulo \(I\), as in Problem~\ref{ec-pr:trace-cong}; \((\Oc/I)(1-k)\) has no invariants under \(I_p\).  Then Problem~\ref{tt-pr:H0}.
\end{description}

\subsection*{\S\ref{sec:relations}}
\begin{description}
\item[\ref{pr:kappa-obeys}] Problem~\ref{ec-pr:trace-cong}, and the upper right entry of \(U(g)\).
\item[\ref{pr:graph-map}] Expand \(U(gh) - 1\) by the cocycle rule.
\item[\ref{pr:graph-AR}] For injectivity, \(\alpha\) kills \((\Lambda/\alpha)^s\).
\item[\ref{pr:graph-toy}] Compute \(\mathcal{M}\) first.
\end{description}

\subsection*{\S\ref{sec:model}}
\begin{description}
\item[\ref{pr:model}] The monomials of weight \(0\), \(+1\) and \(-1\) are \(u^j\), \(zu^j\) and \(yu^j\).  For the cotangent module, keep the linear term of the relation.
\item[\ref{pr:points}] The torus acts by \((y, z) \mapsto (t^{-1}y, tz)\).
\item[\ref{pr:norm}] Elementary divisors: the length is the valuation of the determinant of multiplication by \(a\).  Compute that determinant over \(\overline{\Q}_p\).
\item[\ref{pr:length-bound}] Problem~\ref{pr:norm}.
\end{description}

\begin{thebibliography}{BC}
\bibitem[AM]{AM} M.~F.~Atiyah and I.~G.~Macdonald, \emph{Introduction to Commutative Algebra}, Addison-Wesley, 1969.
\bibitem[BC]{BC} J.~Bella\"iche and G.~Chenevier, \emph{Families of Galois representations and Selmer groups}, Ast\'erisque \textbf{324} (2009).
\bibitem[Ch]{Ch} G.~Chenevier, \emph{The \(p\)-adic analytic space of pseudocharacters of a profinite group and pseudorepresentations over arbitrary rings}, in: Automorphic Forms and Galois Representations, Vol.~1, London Math.\ Soc.\ Lecture Note Ser.\ 414, Cambridge Univ.\ Press, 2014, 221--285; \url{https://arxiv.org/abs/0809.0415}.
\bibitem[GL]{GL} L.~Sernau, \emph{GRH for Dirichlet \(L\)-functions, in Lean~4}, a shorter rewriting of the Lean proof, text by Claude, \url{https://github.com/lsernau/grh-lean}, commit \texttt{72fead1}, readable edition \texttt{grh/}.
\bibitem[MW]{MW} B.~Mazur and A.~Wiles, \emph{Class fields of abelian extensions of \(\mathbf{Q}\)}, Invent.\ Math.\ \textbf{76} (1984), 179--330.
\bibitem[Ro]{Ro} R.~Rouquier, \emph{Caract\'erisation des caract\`eres et pseudo-caract\`eres}, J.~Algebra \textbf{180} (1996), 571--586.
\bibitem[Ta]{Ta} R.~Taylor, \emph{Galois representations associated to Siegel modular forms of low weight}, Duke Math.\ J.\ \textbf{63} (1991), 281--332.
\end{thebibliography}

\end{document}
